\section*{第 \ref{ch:b1:functions} 章 --- 标准函数}

\begin{solution}{exo:b1:functions:1}
$\arcsin\frac{\sqrt 3}{2} = \frac\pi3$;
$\;\arccos\bigl(-\frac12\bigr) = \frac{2\pi}{3}$;
$\;\arctan(-1) = -\frac\pi4$.

$\arcsin\bigl(\sin\frac{5\pi}{6}\bigr)$: $\sin\frac{5\pi}{6} =
\frac12$，且在 $\intcc{-\frac\pi2}{\frac\pi2}$ 内正弦值为
$\frac12$ 的角为 $\frac\pi6$（而非 $\frac{5\pi}{6}$）。

$\arccos\bigl(\cos\frac{7\pi}{4}\bigr)$: $\cos\frac{7\pi}{4} =
\frac{\sqrt 2}{2}$，且在 $\intcc{0}{\pi}$ 内余弦值为该数的角
为 $\frac\pi4$。
\end{solution}

\begin{solution}{exo:b1:functions:2}
$f(x) = \arcsin(2x - 1)$ 要求 $-1 \leq 2x - 1 \leq 1$，即 $x
\in \intcc{0}{1}$。在 $\intoo{0}{1}$ 上，由链式法则及
\cref{prop:b1:functions:arcderiv}可得
\[
f'(x) = \frac{2}{\sqrt{1 - (2x-1)^2}} = \frac{2}{\sqrt{4x - 4x^2}}
= \frac{1}{\sqrt{x(1 - x)}} .
\]
$g(x) = \arctan\sqrt x$ 的定义域为 $x \geq 0$，且对 $x > 0$ 有：
\[
g'(x) = \frac{1}{1 + x} \cdot \frac{1}{2\sqrt x}
= \frac{1}{2\sqrt x\,(1 + x)} .
\]
\end{solution}

\begin{solution}{exo:b1:functions:3}
$\cosh x + \sinh x = \frac{\eu^x + \eu^{-x}}{2} + \frac{\eu^x -
\eu^{-x}}{2} = \eu^x$。由此，对 $n \in \Z$ 有：
\[
(\cosh x + \sinh x)^n = (\eu^x)^n = \eu^{nx}
= \cosh nx + \sinh nx .
\]
\end{solution}

\begin{solution}{exo:b1:functions:4}
设 $y = \arcsin x$，则 $\cos 2y = 1 - 2\sin^2 y = 1 - 2x^2$，故
$\cos(2\arcsin x) = 1 - 2x^2$.

设 $y = \arctan x$，则 $\sin 2y = 2 \sin y \cos y = 2 \tan y \cos^2 y =
\frac{2\tan y}{1 + \tan^2 y}$，故 $\sin(2\arctan x) =
\dfrac{2x}{1 + x^2}$.
\end{solution}

\begin{solution}{exo:b1:functions:5}
由慢到快：
\[
(\ln x)^{1000} \;\ll\; x^{100} \;\ll\; x^{\ln x} \;\ll\;
\eu^{\sqrt x} \;\ll\; \eu^{x/100} .
\]
理由：$(\ln x)^{1000}/x^{100} \to 0$，由
\cref{prop:b1:functions:powerrules}（对数增长慢于幂）。
$x^{100} = \eu^{100\ln x}$，且 $x^{\ln x} = \eu^{(\ln x)^2}$：由于
$(\ln x)^2 - 100\ln x \to +\infty$，后者占优。$x^{\ln x} =
\eu^{(\ln x)^2} \ll \eu^{\sqrt x}$ 因为 $(\ln x)^2/\sqrt x \to 0$
（对数增长慢于幂 $x^{1/4}$，平方后结论依然成立）。最后 $\sqrt x - x/100
\to -\infty$，故 $\eu^{\sqrt x} \ll \eu^{x/100}$。
\end{solution}

\begin{solution}{exo:b1:functions:6}
定义域：$x \neq \pm 1$，共三个区间。在每个区间上，
\[
f'(x) = \frac{\bigl(\frac{2x}{1-x^2}\bigr)'}{1 +
\frac{4x^2}{(1-x^2)^2}}
= \frac{\frac{2(1-x^2) + 4x^2}{(1-x^2)^2}}
{\frac{(1-x^2)^2 + 4x^2}{(1-x^2)^2}}
= \frac{2(1 + x^2)}{(1 + x^2)^2} = \frac{2}{1 + x^2}
= (2\arctan x)' .
\]
故 $f(x) - 2\arctan x$ 在每个区间上为常数。取值：当 $x =
0$ 时，$f(0) = 0$，常数为 $0$，即在 $\intoo{-1}{1}$ 上。当 $x \to
+\infty$，$\frac{2x}{1-x^2} \to 0^-$，故 $f \to 0$，而 $2\arctan x
\to \pi$：常数为 $-\pi$，即在 $\intoo{1}{+\infty}$ 上。由奇函数性质，
则常数为 $+\pi$，即在 $\intoo{-\infty}{-1}$ 上。综上，$f = 2\arctan x$ 在
$\intoo{-1}{1}$ 上，$f = 2\arctan x - \pi$，当 $x > 1$ 时，$f = 2\arctan x +
\pi$，当 $x < -1$ 时。（这正是正切的二倍角公式，
借 $\arctan$ 来理解。）
\end{solution}

\begin{solution}{exo:b1:functions:7}
$\cosh x = 2$：设 $u = \eu^x > 0$，则 $u + u^{-1} = 4$，故 $u^2 - 4u +
1 = 0$，$u = 2 \pm \sqrt 3$：$x = \ln(2 + \sqrt 3)$ 或 $x = \ln(2 -
\sqrt 3) = -\ln(2 + \sqrt 3)$（两解互为相反数，因
$\cosh$ 为偶函数，两者均成立）。等价地，$x = \pm
\operatorname{arcosh} 2$.

$5\cosh x - 4\sinh x = 3$：代入指数形式的定义，
$\frac{5(u + u^{-1}) - 4(u - u^{-1})}{2} = 3$，即 $u + 9u^{-1} =
6$，亦即 $u^2 - 6u + 9 = (u - 3)^2 = 0$：$u = 3$，唯一解为 $x =
\ln 3$.
\end{solution}

\begin{solution}{exo:b1:functions:8}
设 $\alpha = \arctan\frac12 + \arctan\frac13$。由加法公式：
\[
\tan\alpha = \frac{\frac12 + \frac13}{1 - \frac12\cdot\frac13}
= \frac{5/6}{5/6} = 1 .
\]
两个反正切值均位于 $\intoo{0}{\frac\pi4}$ 内（其自变量均在
$\intoo{0}{1}$ 内），故 $\alpha \in \intoo{0}{\frac\pi2}$；该区间内正切值为
$1$ 的角只有 $\frac\pi4$。

马钦公式：设 $\beta = \arctan\frac15$，两次应用二倍角公式：
\[
\tan 2\beta = \frac{2/5}{1 - 1/25} = \frac{5}{12},
\qquad
\tan 4\beta = \frac{2 \cdot 5/12}{1 - 25/144} = \frac{120}{119}.
\]
于是
\[
\tan\Bigl(4\beta - \frac\pi4\Bigr)
= \frac{\frac{120}{119} - 1}{1 + \frac{120}{119}}
= \frac{1/119}{239/119} = \frac{1}{239}.
\]
范围检验：$\beta < \arctan 1 = \frac\pi4$，且事实上 $\tan 4\beta =
\frac{120}{119}$ 接近 $1$，且 $4\beta \in \intoo{0}{\frac\pi2}$
（因 $\beta < \frac\pi8$，且 $\tan\frac\pi8 = \sqrt 2 - 1 > \frac15$）；
故 $4\beta - \frac\pi4 \in \intoo{-\frac\pi4}{\frac\pi4}$，在此区间上
$\arctan$ 是 $\tan$ 的逆：$4\beta - \frac\pi4 =
\arctan\frac{1}{239}$，即马钦公式。
\end{solution}

\begin{solution}{exo:b1:functions:9}
设 $f(x) = x - \sin x$：$f(0) = 0$ 且 $f'(x) = 1 - \cos x \geq 0$，
故 $f \geq 0$ 在 $\R_+$ 上，即 $\sin x \leq x$。

设 $g(x) = \sin x - x + \frac{x^3}{6}$：$g(0) = 0$，$g'(x) = \cos x -
1 + \frac{x^2}{2}$，$g'(0) = 0$，$g''(x) = -\sin x + x = f(x) \geq 0$
在 $\R_+$ 上。于是 $g'$ 单调递增，且 $g'(0) = 0$，故 $g' \geq 0$，
从而 $g$ 单调递增，且 $g(0) = 0$，因此 $g \geq 0$，即 $x -
\frac{x^3}{6} \leq \sin x$ 在 $\R_+$ 上。

因此 $0 \leq \frac{x - \sin x}{x^3} \leq \frac 16$，当 $x >
0$ 时，该比值被夹在 $\frac16$ 的量级附近，
并由\cref{ch:b1:taylor}可知其极限恰为 $\frac 16$。
\end{solution}

\begin{solution}{exo:b1:functions:10}
记 $y = \arcsin x$，则 $2x\sqrt{1 - x^2} = 2\sin y\cos y = \sin
2y$，且 $h(x) = \arcsin(\sin 2y) - 2y$。

当 $\abs x \leq \frac{\sqrt 2}{2}$ 时：$2y \in
\intcc{-\frac\pi2}{\frac\pi2}$，故 $\arcsin(\sin 2y) = 2y$，从而 $h =
0$。

当 $x \in \intoc{\frac{\sqrt 2}{2}}{1}$ 时：$2y \in
\intoc{\frac\pi2}{\pi}$，而
在 $\intcc{-\frac\pi2}{\frac\pi2}$ 中正弦值为 $\sin 2y$ 的角是 $\pi - 2y$：
于是 $h(x) = \pi - 4\arcsin x$。（求导验证：此时
$h'(x) = \frac{2(1 - 2x^2)}{\abs{1 - 2x^2}\sqrt{1 - x^2}} -
\frac{2}{\sqrt{1-x^2}} = \frac{-4}{\sqrt{1 - x^2}}$，即
$-4\arcsin x$ 的导数；又当 $x = 1$ 时，$h(1) = \arcsin 0 - 2\cdot\frac\pi2
= -\pi = \pi - 4\cdot\frac\pi2$。）

当 $x \in \intco{-1}{-\frac{\sqrt 2}{2}}$ 时，由 $h$ 为奇函数，
得 $h(x) = -\pi - 4\arcsin x$。
\end{solution}

\begin{solution}{exo:b1:functions:11}
$g = \arctan \circ \sinh$ 为两个严格递增奇函数的复合，
故其为奇函数且严格递增；当 $x \to
+\infty$ 时，$\sinh x \to +\infty$，故 $g(x) \to \frac\pi2$，且 $g$ 是
到 $\intoo{-\frac\pi2}{\frac\pi2}$ 上的双射（由连续性及
介值定理）。由链式法则并利用
\cref{prop:b1:functions:hyprules}~(1)，得：
\[
g'(x) = \frac{\cosh x}{1 + \sinh^2 x} = \frac{\cosh x}{\cosh^2 x}
= \frac{1}{\cosh x} .
\]
由构造可知 $\tan g(x) = \sinh x$。又因 $g(x) \in
\intoo{-\frac\pi2}{\frac\pi2}$，余弦为正，
\begin{align*}
\cos g(x) &= \frac{1}{\sqrt{1 + \tan^2 g(x)}}
= \frac{1}{\sqrt{1 + \sinh^2 x}} = \frac{1}{\cosh x},\\
\sin g(x) &= \tan g(x) \cos g(x) = \frac{\sinh x}{\cosh x} = \tanh x .
\end{align*}
\end{solution}

\begin{solution}{exo:b1:functions:12}
设 $u = \arctan 2$，$v = \arctan 3$；二者均属于
$\intoo{\frac\pi4}{\frac\pi2}$（其自变量均大于 $1$），故 $u
+ v \in \intoo{\frac\pi2}{\pi}$。由加法公式得
\[
\tan(u + v) = \frac{2 + 3}{1 - 6} = -1 ,
\]
在 $\intoo{\frac\pi2}{\pi}$ 中正切为 $-1$ 的角唯一，
即 $\frac{3\pi}4$，故 $\arctan 2 + \arctan 3 = \frac{3\pi}4$。
（若直接对正切值取 $\arctan$，将得到
$-\frac\pi4$，相差 $\pi$ --- 正是区间定位挽救了
计算，而负分母 $1 - 6$ 恰是
标志着和已越出
$\intoo{-\frac\pi2}{\frac\pi2}$。）加上 $\arctan 1 = \frac\pi4$：
\[
\arctan 1 + \arctan 2 + \arctan 3 = \frac\pi4 + \frac{3\pi}4 = \pi .
\]
\end{solution}

\begin{solution}{pb:b1:functions:1}
\textbf{1.} $f$ 为可导函数之商，其
分母在 $\intoo0{+\infty}$ 上不为零，且
\[
f'(t) = \frac{\frac1t \cdot t - \ln t}{t^2} = \frac{1 - \ln
t}{t^2} ,
\]
当 $t < \eu$ 时为正，在 $t = \eu$ 处为零，当 $t > \eu$ 时为负：
$f$ 在 $\intoc0\eu$ 上严格递增，在
$\intco\eu{+\infty}$ 上严格递减，其最大值为 $f(\eu) = \frac1\eu$。

\textbf{2.} 当 $t \to 0^+$ 时：$\ln t \to -\infty$，且 $\frac1t \to
+\infty$，故 $f(t) \to -\infty$。当 $t \to +\infty$ 时：$f(t) \to 0$
由\cref{prop:b1:functions:powerrules}的增长比较可得
（取 $\beta = \alpha = 1$）。符号与 $\ln t$ 相同，故 $f < 0$ 在
$\intoo01$ 上，$f(1) = 0$，$f > 0$ 在 $\intoo1{+\infty}$ 上。图像
自 $-\infty$ 上升，在 $1$ 处穿过零点，在 $(\eu,
\frac1\eu)$ 处达到峰值，随后衰减至 $0^+$。

\textbf{3.} $f(4) = \frac{\ln 4}4 = \frac{2\ln 2}4 = \frac{\ln 2}2
= f(2)$.

\textbf{4.} 由变化表及介值
定理（高中阶段已使用，在
\cref{ch:b1:continuity}中严格化）。对 $c < 0$：解只存在于 $f < 0$，
即在 $\intoo01$ 上，此处 $f$ 为严格递增双射，
其像为 $\intoo{-\infty}0$，故恰有一解。对 $c = 0$：仅
$t = 1$。对 $0 < c < \frac1\eu$：在 $\intoo1\eu$ 上，$f$ 由
$0$ 增至 $\frac1\eu$，有一解；在 $\intoo\eu{+\infty}$ 上，
$f$ 由 $\frac1\eu$ 减至 $0$，再有一解，共两解。对
$c = \frac1\eu$：仅最大值点 $t = \eu$。对 $c >
\frac1\eu$：无解。

\textbf{5.} 当 $x \neq \eu$ 时，由严格最大性得 $f(x) <
f(\eu)$，即 $\frac{\ln x}x < \frac1\eu$，亦即 $\eu\ln x < x$，
即 $\ln(x^\eu) < \ln(\eu^x)$，故 $x^\eu < \eu^x$。取 $x = \pi$，
则 $\eu^\pi > \pi^\eu$（数值上 $23.14 > 22.46$）。

\textbf{6.} 设 $x, y > 0$，两边均为正，故
\[
x^y = y^x \iff y\ln x = x\ln y \iff \frac{\ln x}x = \frac{\ln y}y
\iff f(x) = f(y) ,
\]
两边除以 $xy > 0$。

\textbf{7.} 设 $f(x) = f(y) = c$，$x \neq y$。若 $c \leq 0$ 或
$c = \frac1\eu$，由第 4 题，方程 $f = c$ 恰有一解，
与题设矛盾。故 $0 < c < \frac1\eu$，且再由第 4 题，
两解中一个属于 $\intoo1\eu$，
另一个属于 $\intoo\eu{+\infty}$：两坐标均大于 $1$，
且二者位于 $\eu$ 两侧。

\textbf{8.} 将 $y = tx$ 代入 $y\ln x = x\ln y$：
\[
tx\ln x = x(\ln t + \ln x)
\iff (t - 1)\ln x = \ln t
\iff \ln x = \frac{\ln t}{t - 1} ,
\]
故 $x = t^{1/(t-1)}$，$y = tx = t^{1 + \frac1{t-1}} =
t^{t/(t-1)}$。反之，对这些值，$\ln y = \frac{t\ln
t}{t-1} = t\ln x$，$y = tx$，故 $\frac{\ln y}y = \frac{t\ln
x}{tx} = \frac{\ln x}x$，即为一个解。每个非平凡解都有
某个比值 $t = y/x \in \intoo0{+\infty} \setminus \{1\}$，故
该参数化是完备的。

\textbf{9.} 取 $t = 2$：$x = 2^{1/1} = 2$，$y = 2^{2/1} = 4$。又
\[
x(1/t) = (1/t)^{\frac1{\frac1t - 1}}
= (t^{-1})^{\frac{t}{1 - t}}
= t^{\frac{t}{t - 1}} = y(t) ,
\]
于是 $y(1/t) = \frac1t\,x(1/t) = \frac{y(t)}t = x(t)$，即
参数变换 $t \mapsto 1/t$ 交换了两个坐标，
这正是方程对称性所要求的。

\textbf{10.} 当 $t \to 1$ 时：$\frac{\ln t}{t-1} \to 1$，
（它是 $\ln$ 在 $1$ 处的差商），故 $x \to \eu^1 = \eu$，且
$y = tx \to \eu$。当 $t \to +\infty$ 时：$\ln x = \frac{\ln t}{t-1}
\to 0$，故 $x \to 1$；而 $\ln y = \frac{t\ln t}{t-1} \to
+\infty$，故 $y \to +\infty$。当 $t \to 0^+$ 时：$\ln x = \frac{\ln
t}{t-1} \to \frac{-\infty}{-1} = +\infty$，故 $x \to +\infty$，
而 $\ln y = \frac{t\ln t}{t-1} \to \frac{0}{-1} = 0$，故 $y \to
1$（用到 $t\ln t \to 0$，见\cref{ex:b1:functions:xx}）。因此该分支
自渐近线 $y = 1$（右侧远处）出发，向上经过
对角线上的 $(\eu, \eu)$，再沿渐近线 $x = 1$ 远去
（上方远处） --- 关于对角线对称，由第 9 题可知。

\textbf{11.} $\ln x(t) = g(t) = \frac{\ln t}{t-1}$，且
\[
g'(t) = \frac{\frac{t-1}t - \ln t}{(t-1)^2}
= \frac{h(t)}{(t-1)^2},
\qquad h(t) = 1 - \frac1t - \ln t .
\]
$h(1) = 0$，且 $h'(t) = \frac1{t^2} - \frac1t = \frac{1 -
t}{t^2} < 0$（$t > 1$）。于是 $h < 0$ 在 $\intoo1{+\infty}$ 上成立，故 $g' <
0$，$x = \eu^g$ 在该区间上严格递减，
由第 10 题，它自 $\eu$ 降至 $1$。

\textbf{12.} 由第 11 题，$t \mapsto x(t)$ 是
$\intoo1{+\infty}$ 到 $\intoo1\eu$ 上的严格递减双射；
记其逆为 $t(x)$（同样严格递减），并令
$\varphi(x) = y(t(x))$。注意 $h < 0$ 在 $\intoo01$ 上也成立（此时
$h' > 0$，$h(1) = 0$），故 $x(\cdot)$ 在
$\intoo01$ 上同样递减；由第 9 题，$y(t) = x(1/t)$ 关于 $t$ 在
$\intoo1{+\infty}$ 上\emph{递增}，由 $\eu$ 升至
$+\infty$。复合得 $\varphi = y \circ t(\cdot)$ 严格递减，
从 $\intoo1\eu$ 映上 $\intoo\eu{+\infty}$，并且
由第 8 题，$x^{\varphi(x)} = \varphi(x)^x$。最后，对公共值
$c = f(x) \in \intoo0{\frac1\eu}$，数对 $\{x,
\varphi(x)\}$ 恰为方程 $f = c$ 的\emph{两元解集}，
（由第 4 题）；将 $\varphi$ 延拓至 $\intoo\eu{+\infty}$ 上，作为其逆映射，则
$\varphi(\varphi(x))$ 回到另一个（即
原来的）元素，即 $\varphi \circ \varphi = \mathrm{id}$。

\textbf{13.} 若 $(x, y)$ 为非平凡整数解，且 $x
< y$，由第 7 题得 $x \in \intoo1\eu$，其中唯一的整数是
$x = 2$。于是 $f(y) = f(2) = \frac{\ln2}2$，且 $y > \eu$；
由第 4 题，方程 $f = \frac{\ln2}2$ 在 $\eu$ 之外恰有一解，
由第 3 题可给出该解：$y = 4$。因此
$(2, 4)$ 及其对换即为全部解。

\textbf{14.} 取 $t = 1 + \frac1n = \frac{n+1}n$，则 $\frac1{t - 1}
= n$，$\frac t{t-1} = n + 1$，故
\[
x_n = \Bigl(\frac{n+1}n\Bigr)^{n},
\qquad
y_n = \Bigl(\frac{n+1}n\Bigr)^{n+1},
\]
显然为有理数，且二者互异（$y_n = t\,x_n \neq x_n$），又 $n =
1$ 时 $x_1 = 2$，$y_1 = 4$。

\textbf{15.} $t = y/x$ 为有理数且 $> 1$；记 $t - 1 =
\frac rs$（既约），故 $t = \frac{s + r}s$，且 $\gcd(s +
r, s) = \gcd(r, s) = 1$，由第 8 题得 $x = t^{s/r}$，从而
\[
x^r = t^s = \frac{(s+r)^s}{s^s} ,
\]
这是一个既约分数（不存在质数同时整除 $s + r$ 与
$s$）。设 $x = \frac pq$ 为既约分数，则 $x^r = \frac{p^r}
{q^r}$ 同样既约，由既约表示的唯一性，
$p^r = (s+r)^s$，$q^r = s^s$。比较
任一质数 $\ell$ 在 $q^r = s^s$ 中的指数：$r \cdot
v_\ell(q) = s \cdot v_\ell(s)$，故 $r$ 整除 $s\,v_\ell(s)$；
由 $\gcd(r, s) = 1$，$r$ 整除 $v_\ell(s)$，
对每个质数 $\ell$ 成立（唯一分解定理，在本阶段不加证明地承认；证明在
\cref{ch:b1:arith}中给出），于是 $s = b^r$，其中 $b$ 为整数。对
$p^r = (s+r)^s$ 作同样的论证，得 $s + r = a^r$。

\textbf{16.} 假设 $a^r - b^r = r$，其中整数 $a > b \geq 1$
且 $r \geq 2$。则 $a \geq b + 1$，由二项式定理
\[
a^r - b^r \geq (b+1)^r - b^r
= \sum_{k=0}^{r-1}\binom rk b^k
\geq \binom r0 + \binom r{r-1} b^{r-1}
\geq 1 + r > r ,
\]
故得矛盾。由第 15 题，$a^r - b^r = (s + r) - s = r$
推出 $r = 1$：$t = 1 + \frac1s$，解为 $(x_s,
y_s)$（第 14 题）。连同互换的数对，
分类即告完成。

\textbf{17.} $f\bigl(\tfrac94\bigr) = \frac{\ln 2.25}{2.25} =
\frac{0.81093}{2.25} = 0.3604$ 与 $f\bigl(\tfrac{27}8\bigr) =
\frac{\ln 3.375}{3.375} = \frac{1.21640}{3.375} = 0.3604$（四位
小数）：相等，正如构造所保证的 --- 于是
$\bigl(\tfrac94\bigr)^{27/8} = \bigl(\tfrac{27}8\bigr)^{9/4}$。

\textbf{18.} 参数 $t_n = 1 + \frac1n$ 严格递减
趋于 $1$。由于 $x(\cdot)$ 在
$\intoo1{+\infty}$ 上严格递减（第 11 题），$x_n = x(t_n)$ 严格
\emph{递增}；由于 $y(\cdot)$ 在该区间上严格递增
（第 12 题），$y_n = y(t_n)$ 严格\emph{递减}。由
第 7 题，$x_n \in \intoo1\eu$ 与 $y_n \in
\intoo\eu{+\infty}$：故 $x_n < \eu < y_n$ 对所有 $n$ 成立。最后
由 $t_n \to 1$ 及第 10 题得 $x_n \to \eu$ 与 $y_n \to \eu$。
$\bigl(1 +
\frac1n\bigr)^n$ 的经典单调收敛性由分支的几何直接得出，
无需新的不等式。

\textbf{19.} 若 $\eu \leq a < b$：$f$ 在
$\intco\eu{+\infty}$ 上严格递减，故 $f(a) > f(b)$，即 $b\ln a > a\ln
b$，从而 $a^b > b^a$。若 $1 < a < b \leq \eu$：$f$ 严格
递增，故 $f(a) < f(b)$，于是 $a^b < b^a$。混合情形：
$2 < \eu < 3$ 且 $2^3 = 8 < 9 = 3^2$；但 $2 < \eu < 5$ 且
$2^5 = 32 > 25 = 5^2$：两种结果均会出现，取决于点
$(a, b)$ 落在分支的哪一侧。

\textbf{20.} 分支与对角线交于 $\eu$ 处，$\varphi$
在该点连续延拓，且 $\varphi(\eu) = \eu$。
对 $\varphi(\varphi(x)) = x$ 在 $x = \eu$ 处求导，由
链式法则得 $\varphi'(\varphi(\eu))\,\varphi'(\eu) =
\varphi'(\eu)^2 = 1$，故 $\varphi'(\eu) = \pm1$；$\varphi$ 为
递减函数，故 $\varphi'(\eu) = -1$。分支与
对角线以斜率 $-1$ 相交，即垂直相交。

\textbf{21.} $y = 3x$ 即 $t = 3$ 的情形：$x = 3^{1/2} = \sqrt3$
且 $y = 3^{3/2} = 3\sqrt3$。验证：$f(\sqrt3) =
\frac{0.54931}{1.73205} = 0.3171$ 与 $f(3\sqrt3) =
\frac{\ln 5.19615}{5.19615} = \frac{1.64792}{5.19615} = 0.3171$：
四位小数下相等，故 $(\sqrt3)^{3\sqrt3} =
(3\sqrt3)^{\sqrt3}$。

\textbf{22.} $n^{1/n} = \eu^{f(n)}$，又 $\exp$ 递增，故
次序由 $f(n)$ 决定：$f(3) = \frac{\ln3}3 \approx
0.3662$ 大于 $f(2) = f(4) = \frac{\ln2}2 \approx 0.3466$（第
3 题），且 $f(5) \approx 0.3219$。最大者为 $\sqrt[3]3$，
相等的一对为 $\sqrt2 = \sqrt[4]4$ --- 即整数对 $(2, 4)$
以另一种形式再次出现。

\textbf{23.} 解集是对角线 $\{(x,
x) : x > 0\}$ 与一条关于对角线对称的单一分支的并，
严格递减，从渐近线 $x = 1$（当 $y \to
+\infty$）趋于渐近线 $y = 1$（当 $x \to +\infty$），
与对角线恰交于一点 $(\eu, \eu)$，该处斜率为 $-1$。
分支上的整点为 $(2, 4)$ 与 $(4, 2)$
--- 仅有这两个 ---，以及有理点 $(x_n, y_n)$，
它们沿分支单调趋向 $(\eu, \eu)$，
却永不到达（$\eu$ 为无理数；分支上的有理点
在一个无理角点处聚集）。

\textbf{24.} (i) 由增长比较得到 $f$ 在
$+\infty$ 处的极限（第 2 题）与 $t\ln t \to 0$（第 10 题），
共同决定了变化表与渐近线。(ii) 中间值定理
经由双射的陈述，
将变化表化为精确的解数
（第 4 题），并化为逆映射 $t(x)$ 的存在性
（第 12 题）。(iii) 唯一分解定理促成了两处
既约形式的确定（第 15 题），即有理点分类的
算术核心。

\textbf{25.} 仅一次求导计算 --- 即 $1 - \ln
t$ 的符号 --- 生成了全部内容：各水平
$c$ 下的解数、分支的形状与渐近线、极值
不等式 $\eu^x \geq x^\eu$、整数解与
有理解的分类，以及 $\bigl(1 +
\frac1n\bigr)^n$ 的单调收敛。这是变化表所体现的思维
经济性：一元研究即可解决二元方程，
因为该方程可经由一个函数分解。
\cref{ch:b1:seq}将重新证明 $(x_n)$ 的收敛性，借助
单调数列的一般理论；\cref{ch:b1:derivative}将
以中值定理严格地建立变化表，而
\cref{ch:b1:taylor}给出变化表无法提供的内容 --- 即
\emph{速度}，即 $x_n \to \eu$ 的收敛（其
阶为 $1/n$）。
\end{solution}
